\section{Applications}\label{sec:applications}
FLOW-DC grew from the practical need to acquire large image collections for biodiversity and agricultural machine learning. These projects connect public metadata, distributed image hosts, dataset curation, and model development. Earlier versions of the acquisition workflow supported the projects below, as reported by the contributing authors. The publications document the resulting datasets and learning tasks. The current \ratio{} and \grad{} rules are evaluated separately in Section~\ref{sec:results}.

\subsection{Insect Identification from Citizen Science}
InsectNet combines image acquisition, representation learning, classification, and deployment for insect identification~\cite{insectnet}. Its data pipeline draws on iNaturalist imagery, with approximately six million labeled images spanning 2,526 insect species. The earlier study, \emph{Deep learning powered real-time identification of insects using citizen science data}~\cite{insect-preprint}, belongs to the same research lineage. These are related stages of one application family, rather than two independent performance evaluations.

The published InsectNet pipeline calls its downloader iNaturalist Scalable Download (iNatSD). The authors identify this acquisition work as part of the lineage leading to FLOW-DC. Its practical requirements---large URL lists, taxonomic metadata, distributed execution, and reusable image output---motivate the interfaces described here.

\subsection{BioTrove and Arboretum}
BioTrove provides a large curated biodiversity image collection and supports models that connect visual and textual information~\cite{biotrove,clip}. The published dataset contains approximately 161.9 million images across 366.6 thousand species. The earlier \emph{Arboretum: A Large Multimodal Dataset Enabling AI for Biodiversity} manuscript~\cite{arboretum} is the first version of the same arXiv record and reports an earlier dataset scale. It is therefore treated as an earlier dataset iteration rather than a separate deployment.

At this scale, acquisition is a workflow problem as much as a transfer problem. Manifest partitions must retain taxonomic and attribution metadata, workers must produce manageable artifacts, and failures must remain visible for later curation. An earlier FLOW-DC workflow supported this acquisition effort. That experience supports the system's emphasis on partitioned scheduling and archive-oriented output, while the current row-level verification strengthens the accountability of subsequent acquisitions.

\subsection{Weakly Supervised Insect Detection}
\emph{Zero-shot insect detection via weak language supervision} studies how image and language information can support insect detection without conventional task-specific supervision~\cite{weak-insects}. The acquisition workflow contributes reusable image collections and metadata that can serve tasks beyond the original classification objective. This use case emphasizes why FLOW-DC preserves the relationship between source rows and delivered images: a dataset assembled once may later support a different sampling rule, annotation strategy, or model family. Downloading is one enabling component; the downstream detection method remains a separate scientific contribution.

\subsection{Weed Identification and Reusable Research Infrastructure}
WeedNet develops a global-to-local approach to weed species identification, with a global model spanning 1,593 species and regional specialization~\cite{weednet}. The work appeared first as a 2025 preprint and was published online in \emph{Nature Communications} in September 2026. The acquisition workflow supported its dataset construction. This application illustrates the need to reuse the same acquisition machinery across biological groups while changing manifest selection and downstream learning procedures.

Across these projects, the recurring requirement is a dependable transition from public image references to organized research data. Research cyberinfrastructure such as CyVerse provides complementary environments for managing and sharing scientific analyses~\cite{cyverse}. FLOW-DC addresses the acquisition portion of that ecosystem. Its value lies in making large transfers easier to organize, adapt, inspect, and reuse, while preserving the source metadata and operational evidence needed by subsequent curation and analysis.
